\section{Introduction}\label{sec:intro}

\paragraph{The missing layer.} Large language models routinely produce \emph{competent} text but rarely produce text whose interpretive multiplicity is deliberately authored. Where humans use a single utterance to gesture toward several incompatible readings simultaneously --- the Wittgensteinian \emph{aspect-shift}~\citep{WittgensteinPI} --- LLMs typically collapse onto a single dominant reading dictated by the prior. We argue that the missing structural element is not scale, alignment data, or instruction-tuning sophistication, but a \emph{recursive self-reflexivity layer}: an explicit operator that re-examines the model's own output against the same constraint manifold that produced it, is permitted to revise the surface when the alternatives are not yet visible in it, and gates that revision on a measurable free-energy signal rather than on the prompt-template's exhortation to ``check your work''. The contemporary literature on iterative self-refinement~\citep{MadaanEtAl2023SelfRefineNeurIPS,ShinnEtAl2023Reflexion,BaiEtAl2022ConstitutionalAI} is converging on the same insight from the engineering side --- giving an LM a critique step buys real quality --- but it leaves open the gating question: \emph{when} should the model revise? On every item, regardless of cost? Only when an internal signal says the draft is suspect? Only when a learned policy says revision will pay off?

\paragraph{The Pratyabhij\~n\=a tradition has a precise answer.} Abhinavagupta's Pratyabhij\~n\=a (``recognition'') school of Kashmir \'Saiva phenomenology developed an unusually fine-grained taxonomy of the layer above first-pass cognition~\citep{Lawrence1996TantricArgumentPEW,abhinavagupta-isvarapratyabhijna,pratyabhijnahrdayam-singh}. The five-\iast{\'sakti} cascade (\iast{cit}, \iast{\=ananda}, \iast{icch\=a}, \iast{j\~n\=ana}, \iast{kriy\=a}) describes the trajectory by which an undifferentiated luminous ground (\iast{prak\=a\'sa}) becomes a determinate utterance, and the \iast{vimar\'sa} layer is the moment of self-touching --- the recognition that the utterance was authored by the very same ground that contemplates it. \iast{Apohana-\'sakti} (the contrastive exclusion of all-other-than-this) is precisely the geometric move that Bayesian model reduction~\citep{Friston2018BayesianModelReductionArxiv,FristonPenny2011,FristonEtAl2016BayesianModelReductionEmpiricalBayesDCM} formalises as posterior switch under a parameter-pruning prior, and that the broader free-energy programme~\citep{Friston2010FreeEnergyUnifiedBrainTheory,Friston2017ActiveInferenceCuriosity,Friston2013LifeAsWeKnowItJRSAInterface,Heins2022pymdp} treats as the operational substrate of insight.

The Pratyabhij\~n\=a tradition's contribution to a recursive-LLM architecture is not a metaphor. It is a precise structural answer to two questions the contemporary literature is not yet asking sharply: (a) the gating question above (the answer is: gate on \iast{vimar\'sa-\'sakti}, the recursive self-touching that occurs when the system recognises its own \iast{prak\=a\'sa} as the source of its current draft), and (b) the \emph{posterior-pruning} question (the answer is: prune via \iast{apohana}, contrastive exclusion against the \emph{not-this} set, formally equivalent to BMR with a sparsifying prior). This paper does not claim the Pratyabhij\~n\=a tradition is true in any metaphysical sense; it claims the architecture is correct, falsifiable when implemented, and useful as a design source for the gating and pruning layers in modern LLM cascades.

\paragraph{What v0.4 reports.} This paper is v0.4 of the \emph{Pratyabhij\~n\=a Creative Engine} (PCE), the executable instantiation of the architecture above. PCE is implemented as a Cursor / Claude Code plugin and as a standalone \texttt{python -m pce} CLI; both share a single substrate (the \texttt{claude --print} OAuth-CLI binary against either an Anthropic API key or AWS Bedrock). v0.4 is deliberately framed as a \emph{mechanism study}: at the available pilot $n$ (4 domains, 25 paired items in total) the headline cascade-vs-bare contrasts are statistically inconclusive, but several mechanism-level signals are robust enough to be load-bearing for the architecture's next iteration. The contributions are:

\begin{enumerate}[leftmargin=2em,itemsep=2pt,topsep=2pt]
\item \textbf{A four-arm Bedrock pilot} (Haiku as cascade scorer, Sonnet-4 as judge) under a bytewise-isolated CLI substrate, addressing the v0.2 reviewer's contamination critique. The merged ledger reports {V04\_COST\_TOTAL\_USD}~over {V04\_COST\_N\_CALLS}~Bedrock calls. Per-domain Hedges' $g$ (H1.v4 through H4.v4) ranges $[-0.32,\,+0.32]$ with no individual domain crossing the BCa zero line; the H5.v4 fixed-effects pool is $g={V04\_H5\_POOLED\_G}$, 95\% CI $[{V04\_H5\_CI\_LO},\,{V04\_H5\_CI\_HI}]$, $n={V04\_N\_PAIRED\_TOTAL}$.
\item \textbf{A within-cascade revision-quality test} (H8a, $n={V04\_H8A\_N}$): the shadow revision strictly beats the draft on average, $g={V04\_H8A\_G}$, $p\approx{V04\_H8A\_P}$. The cascade reliably \emph{produces} a better surface; the open question is whether it reliably \emph{commits} the better surface.
\item \textbf{A four-policy commit multiplexer} (H8c) that ranks \texttt{always\_revise} above \texttt{learned\_gate} above \texttt{event\_gated} above \texttt{always\_draft} on paired delta vs.\ \texttt{haiku\_bare}; the headline is {V04\_H8C\_WINNER}. The \iast{vimar\'sa} event gate as currently defined under-fires (H8b: event-gated F1 $={V04\_H8B\_EVENT\_F1}$ vs.\ a learned gate's F1 $={V04\_H8B\_LEARNED\_F1}$).
\item \textbf{A judge-vs-scorer agreement test} (H9, $n={V04\_H9\_N}$): the Sonnet judge and the Haiku composite-score head are not measuring the same thing on the cascade items (Spearman $\rho={V04\_H9\_RHO}$, sign-agreement {V04\_H9\_SIGN\_AGREEMENT}). We treat this not as a defect but as a load-bearing finding for any cascade architecture that uses a small-LM scorer to gate a large-LM commit.
\end{enumerate}

\paragraph{The unmerged-state question.} An honest reader of this repository will notice that the Phase 7 results that drive every number above were left for several days on the \texttt{v0.4-mechanism-study} branch, never merged to \texttt{main}. We address this directly. The v0.3 SPEC closed with a ``negative-result obligation'': any pilot that failed to find the predicted directional effect must be reported \emph{as that finding}, not framed as a methodological caveat. The v0.3 prose acknowledged the obligation but did not enact it; the prose framed the inconclusive headline as a sample-size limitation and set up additional control arms for a v0.4 rerun, which is what Phase 7 then ran. v0.4 inherits the obligation and discharges it: the headline cascade-vs-bare result is inconclusive at $n=25$ and we say so in the abstract; the load-bearing mechanism findings (H8a, H8c, H9, H8b) are foregrounded for what they actually are; the speculative scope-cuts (chandas-aware scorer, Hopfield \iast{\=alayavij\~n\=ana} substrate, learned gate trained on more data) are deferred to v0.5 with explicit falsification criteria. The unmerged-state evidence-then-prose discipline is, in retrospect, a feature of the pipeline, not a bug: the artefacts on the branch are the truth, and v0.4 is the prose rebuild around them.

\paragraph{Organisation.} \S\ref{sec:related} situates PCE among recent work on Indic-philosophical AI architectures, BMR, LLM-as-judge methodology, iterative self-refinement, and computational Sanskrit. \S\ref{sec:pratyabhijna_background} specifies the Pratyabhij\~n\=a primitives. \S\ref{sec:active_inference_background} states the active-inference and BMR primitives, including the honest scope of free-energy in PCE (best-of-K composite scoring, not exact variational inference). \S\ref{sec:engine} defines the operators and the consolidation cycle. \S\ref{sec:plugin} describes the plugin surface. \S\ref{sec:methods} formalises the v0.4 experimental protocol, including the Bedrock substrate and the commit-policy multiplex. \S\ref{sec:substrate} documents substrate isolation and portability across Cursor, Claude Code, and the standalone CLI. \S\ref{sec:benchmarks} describes items and scoring. \S\ref{sec:results} reports the five hypothesis families (H1--H4, H5, H8a, H8b/c, H9). \S\ref{sec:discussion} interprets the findings across eight subsections. \S\ref{sec:honest_ai_claims} is a dedicated section on what the cascade can and cannot claim about ``creativity'' on the available evidence. \S\ref{sec:showcase} walks the nine release-quality showcase outputs, each with its full draft / shadow-revision / commit trace. \S\ref{sec:limitations} and \S\ref{sec:conclusion} close.
